%==================================================================

Вместо того чтобы каждый раз проверять аксиомы для $\et(\Shv{F})$,
можно <<прокачать>> предпучок внутри заранее готового пучка.

\begin{ass}
Пусть $\Shv{H}$ --- пучок на $X$, а $\Shv{F}\subset\Shv{H}$ ---
подпредпучок: семейства $\Shv{F}(U)\subseteq\Shv{H}(U)$,
согласованные с ограничениями $\rho$.
\end{ass}

\begin{defn}[секции $\Shv{F}^{+}$]
Для открытого $U\subset X$ положим
\begin{equation}\label{eq:Fplus}
 \Shv{F}^{+}(U)=\bigl\{\,s\in\Shv{H}(U)\ \bigm|\ \exists\ \text{покрытие }
 U=\textstyle\bigcup_i U_i\ \text{ и }\ \res^{U_i}_{U}(s)\in\Shv{F}(U_i)
 \ \forall i\bigr\},
\end{equation}
то есть $\Shv{F}^{+}(U)$ --- секции пучка $\Shv{H}$, которые
\emph{локально} являются секциями предпучка $\Shv{F}$.
\end{defn}

\begin{thm}\label{thm:Fplus}
$\Shv{F}^{+}$ --- пучок, и $\Shv{F}^{+}\simeq\Shv{F}^{\#}$.
\end{thm}

Сразу видно, что $\Shv{F}^{+}$ --- пучок: единственность и склейка
наследуются от $\Shv{H}$ (покрытие для двух секций можно объединить,
а для склейки --- взять объединение двух покрытий). Вся работа ---
в изоморфизме.

\subsection{План доказательства}
Вместо прямой проверки используется теорема~\ref{thm:hom} и
Факт~\ref{fact:hom}:
\begin{conditions}
\item построить естественный изоморфизм функторов
 $\Hom(\Shv{F},-)\simeq\Hom(\Shv{F}^{+},-)$ на пучках;
\item изоморфизм $\Hom(\Shv{F},-)\simeq\Hom(\Shv{F}^{\#},-)$ уже есть
 (теорема~\ref{thm:hom});
\item тогда $\Hom(\Shv{F}^{+},-)\simeq\Hom(\Shv{F}^{\#},-)$ естественно
 (композиция изоморфизмов), и <<Факт>> даёт
 $\Shv{F}^{+}\simeq\Shv{F}^{\#}$.
\end{conditions}

\subsection{Построение $\eta_{\Shv{G}}$}
Для пучка $\Shv{G}$ и $\varphi\colon\Shv{F}\to\Shv{G}$ отображение
\[
 \eta_{\Shv{G}}(\varphi)=\bar\varphi\colon\Shv{F}^{+}\to\Shv{G}
\]
строится так.
\begin{conditions}
\item Пусть $s\in\Shv{F}^{+}(U)$. По определению~\eqref{eq:Fplus} есть
 покрытие $U=\bigcup_iU_i$ и $s_i:=\res^{U_i}_{U}(s)\in\Shv{F}(U_i)$.
\item Положим $t_i:=\varphi_{U_i}(s_i)\in\Shv{G}(U_i)$.
\item Семейство $(t_i)$ согласовано: на $U_i\cap U_j$ равенство
 $s_i|=s_j|$ выполняется (оба --- ограничения одного $s$), а $\varphi$
 --- гомоморфизм, поэтому $\varphi(s_i)|=\varphi(s_j)|$.
\item Аксиома (F2) \emph{пучка $\Shv{G}$} даёт единственное
 $t\in\Shv{G}(U)$ с $t|_{U_i}=t_i$. Полагаем $\bar\varphi(s):=t$.
\item Независимость от выбора покрытия: для другого покрытия полученные
 локальные значения на пересечениях совпадают, и снова работает (F2)
 (либо (F1) --- для единственности).
\end{conditions}

Ясно, что $\bar\varphi\circ i=\varphi$ (берём покрытие из одного
множества $U$). Диаграмма~\ref{dgm:eta} фиксирует, что получилось:
$\eta_{\Shv{G}}$ --- естественное по $\Shv{G}$ преобразование
$\Hom(\Shv{F},-)\Rightarrow\Hom(\Shv{F}^{+},-)$.

\begin{diagram}{Слева: $\eta_{\Shv{G}}$ --- естественное по $\Shv{G}$
 преобразование $\Hom(\Shv{F},-)\Rightarrow\Hom(\Shv{F}^{+},-)$;
 справа: $\bar\varphi$ --- единственное продолжение $\varphi$
 до $\Shv{F}^{+}$.\label{dgm:eta}}
\begin{tikzcd}[cd style]
 \Hom(\Shv{F},\Shv{G}_1) \ar[r,"\alpha\circ-"] \ar[d,"\eta_{\Shv{G}_1}"'] &
 \Hom(\Shv{F},\Shv{G}_2) \ar[d,"\eta_{\Shv{G}_2}"] \\
 \Hom(\Shv{F}^{+},\Shv{G}_1) \ar[r,"\alpha\circ-"'] &
 \Hom(\Shv{F}^{+},\Shv{G}_2)
\end{tikzcd}
\qquad
\begin{tikzcd}[cd style,column sep=huge,row sep=huge]
 \Shv{F} \ar[r,"i"] \ar[dr,"\varphi"'] & \Shv{F}^{+} \ar[d,"\bar\varphi"] \\
  & \Shv{G}
\end{tikzcd}
\end{diagram}

Пояснение к диаграмме~\ref{dgm:eta}. Слева --- квадрат
натуральности $\eta_{\Shv{G}}$: домножение морфизма предпучков
на $\alpha\colon\Shv{G}_1\to\Shv{G}_2$ и переход от
$\Hom(\Shv{F},-)$ к $\Hom(\Shv{F}^{+},-)$ коммутируют, т.\,е.\
$\eta$ --- преобразование функторов. Справа --- из чего состоит
компонента $\eta_{\Shv{G}}$: к морфизм $\varphi\colon\Shv{F}\to
\Shv{G}$ единственно продолжается до $\bar\varphi\colon\Shv{F}^{+}
\to\Shv{G}$ с $\bar\varphi\circ i=\varphi$ (треугольник), и
именно $\eta_{\Shv{G}}(\varphi)=\bar\varphi$.

\subsection{Почему $\eta_{\Shv{G}}$ --- изоморфизм}
Остаётся проверить, что каждая компонента $\eta_{\Shv{G}}$ ---
биекция; именно это позволяет применить Факт~\ref{fact:hom}.
\begin{bullets}
\item \emph{Инъективность}: пусть $\eta_{\Shv{G}}(\varphi_1)=
 \eta_{\Shv{G}}(\varphi_2)$. Продолжения совпадают, поэтому после
 композиции с каноническим $i\colon\Shv{F}\to\Shv{F}^{+}$ получаем
 $\varphi_1=\varphi_2$. Иными словами, отображение ограничения
 $\Hom(\Shv{F}^{+},\Shv{G})\to\Hom(\Shv{F},\Shv{G})$ является левым
 обратным к $\eta_{\Shv{G}}$.
\item \emph{Сюръективность}: пусть $\psi\colon\Shv{F}^{+}\to\Shv{G}$.
 Рассмотрим продолжение $\overline{\psi\circ i}$. Для каждого
 $s\in\Shv{F}^{+}(U)$ существует покрытие $U=\bigcup U_i$, на котором
 $s|_{U_i}=i(s_i)$ для $s_i\in\Shv{F}(U_i)$. На каждом $U_i$ оба
 сечения $\overline{\psi\circ i}(s)$ и $\psi(s)$ равны
 $\psi(i(s_i))$. По (F1) для пучка $\Shv{G}$ они равны на $U$.
 Значит, $\eta_{\Shv{G}}(\psi\circ i)=\psi$.
\end{bullets}

Итак $\Hom(\Shv{F},-)\simeq\Hom(\Shv{F}^{+},-)$ естественно, и
Факт~\ref{fact:hom} даёт $\Shv{F}^{+}\simeq\Shv{F}^{\#}$.
\hfill$\square$

\begin{bookbox}
\booksrc{Манин, \S\,2.1.6---2.1.7, с.~112---113: <<секции $\Shv{P}^{+}(U)$
 определяются как некоторые наборы ростков $\prod_{x\in U}\Shv{P}_x$,
 которые в естественном смысле согласованы>>; там же описаны \emph{два}
 процесса <<прокачивания>>: (1) \emph{уменьшить} число сечений,
 отождествляя те, что начинают совпадать на мелком покрытии; (2)
 \emph{увеличить} число сечений, добавляя склеенные из согласованных
 наборов. Теорема-определение 2.1.7: $\Shv{P}^{+}$ --- пучок и
 $(\Shv{P}^{+})_x=\Shv{P}_x$.
\par\smallskip
 Годеман, Замечание 1.2.3, с.~133 --- пучок, порождённый предпучком
 ($\Shv{F}^{\#}=\tilde{\Shv{F}}$); Годеман, п.~3.9, с.~183 ---
 <<локальное равенство» в $\Shv{F}(X)$.
}
\end{bookbox}

%==================================================================
